\documentclass{article}
\usepackage{amsmath}
\usepackage{amssymb}
\usepackage{hyperref}
\usepackage{qcircuit}

\title{Lecture Notes on Shor's Algorithm}
\author{MHJ}
\date{\today}

\begin{document}

\maketitle

\tableofcontents
\newpage

\section{Introduction to Shor's Algorithm}

\subsection{Historical Context}
Shor's algorithm, introduced by Peter Shor in 1994, revolutionized the field of cryptography. It demonstrated that a quantum computer could solve integer factorization and the discrete logarithm problem efficiently, posing a significant threat to classical cryptosystems such as RSA.

\subsection{Problem Statement}
The objective of Shor's algorithm is to factorize a large composite integer \( N \) into its prime factors. Given \( N = p \times q \), where \( p \) and \( q \) are unknown distinct primes, the problem of factorization can be reduced to finding an integer \( a \) such that \( 1 < a < N \) and \( \gcd(a, N) = 1 \).

\section{Classical Component}

\subsection{Reduction to Order-Finding}
The key reduction in Shor's algorithm lies in transforming factorization into an order-finding problem:
\begin{itemize}
    \item Choose a random integer \( a \) such that \( 1 < a < N \) and \( \gcd(a, N) = 1 \).
    \item Determine the smallest positive integer \( r \) such that \( a^r \equiv 1 \pmod{N} \).
\end{itemize}
The integer \( r \) is known as the \textit{order} of \( a \) modulo \( N \), and it helps in discovering the factors of \( N \).

\subsection{Continued Fraction Expansion}
After obtaining the rational approximation from the quantum algorithm, classical computation using continued fraction expansion aids in deducing the correct order \( r \).

\section{Quantum Component}

\subsection{Quantum Fourier Transform (QFT)}
The Quantum Fourier Transform is crucial to Shor's algorithm. It leverages quantum parallelism and interference to efficiently estimate the periodicity of a function.

\[ \text{QFT}: \ket{x} \mapsto \frac{1}{\sqrt{2^n}} \sum_{y=0}^{2^n-1} e^{2\pi ixy/2^n} \ket{y} \]

\subsection{Quantum Order Finding Procedure}
The quantum part of Shor's algorithm is designed to find the period \( r \) of the modular exponentiation function \( f(x) = a^x \bmod N \):

\[ \Qcircuit @C=1em @R=1em {
    \lstick{\ket{0^{n}}} & \qw & \gate{\mathrm{H^{\otimes n}}} & \qw & \multigate{1}{U_f} & \qw & \measureD{} \\
    \lstick{\ket{0^{n}}} & \gate{\mathrm{H^{\otimes m}}} & \qw & \qw & \ghost{U_f} & \gate{\mathrm{QFT}^{-1}} & \measureD{}
}\]

\subsubsection{Details of the Circuit}
\begin{itemize}
    \item Start with two registers. Apply Hadamard gates to obtain superposition.
    \item Implement the unitary operation \( U_f \) for modular exponentiation \( f(x) = a^x \bmod N \).
    \item Perform an inverse QFT to extract the order \( r \) from the phase information.
\end{itemize}

\section{Putting It All Together}

\subsection{Algorithm Steps}
1. Choose a random integer \( a \).
2. Use a quantum circuit to find the order \( r \).
3. If \( r \) is even and \( a^{r/2} \not\equiv -1 \pmod{N} \), compute \( \gcd(a^{r/2} \pm 1, N) \) to obtain non-trivial factors.
4. If unsuccessful, repeat with new \( a \).

\section{Error Analysis and Success Probability}

\subsection{Probability of Success}
The probability that a randomly chosen \( a \) leads to successful factoring is generally greater than 50\%. The quantum component requires carefully managed resources for high precision.

\subsection{Error Sources}
Error correction is crucial, as any errors in operations or QFT can significantly impact the accuracy of period estimation and the eventual factorization.

\section{Conclusion}

Shor's algorithm is an exemplary quantum algorithm demonstrating exponential speedup over classical counterparts. It ignites both excitement for potential quantum advancements and a reconsideration of current cryptographic standards.

\end{document}
